% !TeX root = ../Harmonia-EA-Capability-Map.tex
% Copyright (c) 2026 Mark Hunter. SPDX-License-Identifier: GPL-3.0-or-later
% Generated from the explicit 0.2-draft model; all assertion status is retained.
\chapter{Reading the proposed architecture}\label{model:guide}\label{guide}
Version 0.2-draft is Step 2: a coherent proposed architecture for semantic and structural review. It does not replace the governed Markdown RAD or declare any architecture domain complete. The 0.1 publication remains independently reproducible in the parent directory. Its identities, definitions, semantic boundaries, source traceability and established relationships are preserved.

Element, naming, direct-placement, complete-path and relationship status are separate assertions. A Canonical responsibility can have a Proposed CT3 placement and Proposed structural identifier. A Candidate concern remains Candidate even when a best placement is proposed. Source uncertainty is retained in the governed appendix; the proposal makes a current judgement rather than leaving the concept floating.

Capability, Feature, Function, Service and Process remain different constructs. Every Feature is framed by BEC CT3 with complete CT1/CT2 ancestry. Directly Capability-owned behaviour remains valid. Services expose a client contract; Processes chain Functions where meaningful. None of those semantics is a mandatory Feature-to-Function-to-Service-to-Process hierarchy.

The fourteen Canonical EC responsibilities become CT2s under three Proposed CT1 frames. The application of the structural ID grammar to ECs is itself Proposed; EC-01 through EC-14 remain stable Reference IDs and historical structural aliases. Reuse does not mandate centralised services, separate processes or distributed responsibility.

Downstream Application Component derivation may use this current proposed model as its working assumption, while retaining every assertion status. Derive cohesive responsibility centres from capabilities; do not force them to reproduce current Mneme, Mnemosyne, Ponos, Calliope or other component boundaries. No capability in this publication allocates an application component, module, process or deployment.

AX-01/04 preserve healthcare meaning over machinery; AX-05 separates active use from durable establishment; AX-06 qualifies source authority; AX-07/08/09 distinguish governed action, meaningful evidence and diagnostics; AX-10/11/12 preserve operational integrity; AX-13 governs boundary publication; AX-14/15/17 preserve distinctions, uncertainty and authority; AX-16/18 preserve entity progression without inferred domain authority. CM-R01 through CM-R05 are reviewed through complete placement, bounded responsibility, explicit definitions, cohesive roots and unique current display names.

\begingroup\small\setlength{\TableWidth}{\dimexpr\linewidth-4\tabcolsep-3\arrayrulewidth\relax}
\begin{longtable}{|>{\raggedright\arraybackslash}p{0.7\TableWidth}|>{\raggedright\arraybackslash}p{0.3\TableWidth}|}
\hline\rowcolor{TableHead}\textbf{Standing} & \textbf{Elements}\\\hline\endfirsthead
\hline\rowcolor{TableHead}\textbf{Standing} & \textbf{Elements}\\\hline\endhead
\hline\endfoot
Canonical & 716\\\hline
Candidate & 3\\\hline
Proposed & 143\\\hline
\end{longtable}\endgroup
\subsection*{Material decisions to review}

\textbf{HIGH — BEC decomposition.} Seven cohesive proposed roots replace dependence on the five contextual views. Entity administration remains CT2 with distinct CT3 identity/\allowbreak{}relationship/\allowbreak{}definition abilities. Most care-administration, care-context and operational responsibilities are themselves CT3s beneath newly proposed CT2 frames. This avoids giving every existing heading a one-child intermediate tier.


\textbf{HIGH — Integrated Enterprise hierarchy.} Three proposed roots organise the fourteen Canonical responsibilities as CT2s. Selected EC-04, EC-06, EC-07 and EC-10 responsibilities are decomposed where materially different reusable abilities exist. Other ECs deliberately stop at CT2. Information Authority Qualification is a new Proposed CT2; it represents credibility standing, not clinical truth or independent service assurance.


\textbf{HIGH — Collective service readiness.} The existing Candidate readiness concern receives Proposed BC CT3 placement beneath a new preparedness CT2 in Care Access \& Coordination. A distinct Proposed BEC assembles dependency evidence. The alternative placement in Health Service Direction \& Stewardship remains plausible; both the business placement and information-versus-decision boundary need review.


\textbf{HIGH — Responsibility does not follow use.} Clinical facts, professional privilege, source roster/\allowbreak{}booking authority, public-health decisions, research approval and governance approval remain with their mandated owners. Audit capture, lineage, operational recovery and assurance conclusions retain different responsibilities. No application component is allocated by these proposals.


\textbf{HIGH — Clinical contexts overlap.} The proposed longitudinal/\allowbreak{}community, facility/\allowbreak{}episodic and specialist CT2 frames organise integration purposes. They are not exclusive classes of care. Acute capture, inpatient clinical progression and screening recall retain their canonical composition-level sufficiency decisions; no specialised clinical capture engines are inferred.


\textbf{MEDIUM — Feature completeness.} Additional Features frame specific missing behaviour: preliminary/\allowbreak{}corrected diagnostic information, order outcomes, historical reconstruction, definition/\allowbreak{}binding/\allowbreak{}profile governance, local capacity, governed reference-content revision, binary content, entity-relative activity, administrative recovery and independent assurance. Existing directly owned Functions remain Capability-owned; proposed realisation links do not silently change accountability.


\textbf{MEDIUM — Name and qualification boundaries.} The current display name Health Information Exchange Enablement is Proposed for the Canonical BEC. Its stable ID, label, established name and governed definition remain unchanged. This resolves the current-model name collision without claiming that RADS has accepted the rename. Clinical Qualification Management and physical badge/\allowbreak{}access Credential Management remain different capabilities; the inherited Operational Credential Check Feature retains its clinical qualification definition despite its potentially confusing name.


\textbf{MEDIUM — Prior candidate reconciliation.} The old Enterprise coexistence branch is removed. Generic Context Establishment, State Monitoring, Exception Management and retrieval patterns are retained as scoped concerns/\allowbreak{}reuse evidence where they describe no additional Feature outcome. Useful specific deltas receive explicit proposals; no candidate table adjacency becomes an established behavioural edge.


\textbf{LOW — Structural conventions.} New semantic local codes are category-qualified. The proposed EC hierarchy uses the existing typed structural grammar as an explicit Step 2 extension. Historical EC-01 through EC-14 structural tokens resolve directly to their current identities; they are not reallocated to another element.

\subsection*{Preserved tensions and limitations}

The v0.1 source appendix deliberately still says some tiers and derivations are unresolved. Those are truthful baseline statements; this publication proposes current answers without changing that history. Clinical device maintenance remains Candidate despite proposed membership beside the canonical seven-child set. Sharing-programme meaning is weakly evidenced and receives a low-confidence Proposed elaboration rather than a fabricated canonical definition. The governance publication BEC is a narrow proposal for BC-17 support, whose RAD relevance and enablement questions remain unadjudicated.


The Strategy assurance derivation retains earlier statements that downstream Functions/\allowbreak{}Services/\allowbreak{}Processes remain unestablished, while the current Domain03 assurance catalogue explicitly establishes them. Both sources are preserved; the five established Functions and existing exposure/\allowbreak{}progression relationships form the semantic anchor here. Similarly, the bounded B2.1 account records a later mapping state than some Domain05 overview paragraphs. No programme completion or authority precedence is inferred merely from those inconsistent status summaries.


This is a complete current proposal within Harmonia's bounded capability scope, not an exhaustive hospital enterprise architecture. Externally provided authentication, roster/\allowbreak{}booking machinery, clinical decision-making and uncommitted R3 master-person reconciliation remain boundary context. Detailed assurance information models, universal lifecycle states, outcome taxonomies, binary custody technology and physical deployment remain downstream questions.
